\chapter{2018 年试题（当年科目代码 819，旧大纲）}

\begin{tishi}
\textbf{关于科目代码：}2018 年西华大学该科目的\textbf{考试科目代码为 819}（名称为"工程流体力学"），
与现行的 818 不同。819 属\textbf{旧大纲}，命题范围比现行 818 大纲宽，
\textbf{含明渠流与堰流（本套计算题第 11 题）等超纲内容}，复习时按最低重要度处理即可。

\textbf{试卷说明（原卷原文）：}试题附在答卷内交回，答案必须写在答题纸上，写在试题纸上无效。

\textbf{结构：}选择题 15 分（5 小题，每小题 3 分）、判断题 20 分（10 小题，每小题 2 分）、
填空题 30 分（8 小题，注明每空 3 分）、计算题 85 分（1--7 题必做，8、9、10、11 四题
\textbf{任选做两道}）。满分 150 分。

\textbf{注意：}判断题与填空题的\textbf{大题顺序原卷为"二、判断题"在前、"三、填空题"在后}，
此处照原卷顺序编排。
\end{tishi}

\section{一、选择题（共 15 分，每小题 3 分）}

\begin{ti}
  \item 作用于流体的质量力包括（\quad）。\hfill\xstarfull{3}\yiju{E4 2018 选择 1（作用于流体的质量力）；E5 题库选择 3（体积力有势）}\nandu{易}
  \fourchoices{压力}{摩擦力}{重力}{表面张力}

  \item 渐变流是（\quad）。\hfill\xstarfull{3}\yiju{E1 slide16 简答（均匀流的特点）}\nandu{易}
  \fourchoices{过流断面上速度均匀分布的流动}{流线近于平行直线的流动}%
  {速度不随时间变化的流动}{沿程各断面流量相同的流动}

  \item 如图 1 所示的封闭容器上装有 U 型水银测压计，其中 1、2、3 点位于同一水平面上，
  其压强关系为（\quad）。\hfill\xstarfull{4}\yiju{E4 2018 选择 3（U 形水银测压计）；E4 2015 填空（U 形压差计）}\nandu{中}
  \fourchoices{$p_1>p_2>p_3$}{$p_1=p_2=p_3$}{$p_1<p_2<p_3$}{$p_2<p_1<p_3$}
  \tuyi{原卷图 1：一封闭容器侧壁接一根 U 型水银测压计；管中 1、2、3 三点位于同一水平面上，
  1 点在容器内的水中（水银面以上），2 点在水与汞的分界面处，3 点在水银面以上的水柱中、
  与 2 点处于同一高程。原卷未给出各段液柱高度数值。（原卷影印见本册附录）}

  \item 流量为 $Q$、流速为 $v$ 的射流冲击一块与流向垂直的平板，则平板受到的冲击力为（\quad）。
  \hfill\xstarfull{5}\yiju{E4 2015 计算（水平弯管求 $F_x,F_y$）；E4 2026 计算 3（教材习题 4.12）；E1 slide16}\nandu{中}
  \fourchoices{$Qv$}{$Qv^{2}$}{$\rho Qv$}{$\rho Qv^{2}$}

  \item 长度 $l$、速度 $v$、重力加速度 $g$ 的无量纲组合是（\quad）。\hfill\xstarfull{4}\yiju{E1 slide33（量纲转化）；E2 slide5}\nandu{中}
  \begin{choices}
    \item $\dfrac{v}{gl}$
    \item $\dfrac{lv}{g}$
    \item $\dfrac{gl}{v}$
    \item $\dfrac{v^{2}}{gl}$
  \end{choices}
\end{ti}

\begin{tishi}
\textbf{第 4 题：}原卷各项中流量与流速的乘积符号 OCR 不清（疑有系数 $\rho$ 与平方号的位置差异），
按四个选项的字面转写，请对照原卷影印核实。

\textbf{第 5 题：}原卷四个选项按分式竖排若干行，OCR 仅拾得分子与分母而丢失了对应关系；
此处按原卷出现的四个式子转写，\textbf{选项与分数的对应顺序请对照原卷影印核实}。
\end{tishi}

\section{二、判断题（正确的打"√"，错误的打"×"，每小题 2 分，共计 20 分）}

\begin{ti}
  \item 伯努利方程中 $z+\dfrac{p}{\rho g}+\dfrac{v^{2}}{2g}$ 表示单位质量流体具有的机械能。
  \hfill\xstarfull{5}\yiju{E1 slide16"计算题：伯努利方程"；E1 slide33 与 E2 slide4"第四章＝考试计算题重点"}\nandu{易}

  \item 液体不能承受拉应力。\hfill\xstarfull{2}\yiju{概念题素材；E6 教材 1.1}\nandu{易}

  \item 静水总压力的压力中心就是受压面的形心。\hfill\xstarfull{5}\yiju{E4 2015 计算（圆柱闸门）；E1 slide33}\nandu{中}

  \item 恒定流动中，流线和迹线总是重合的。\hfill\xstarfull{5}\yiju{E4 2015 计算 1（流线方程、流动方向判断）}\nandu{中}

  \item 圆管湍流水力光滑区的沿程阻力系数 $\lambda$ 和雷诺数 $\Rey$ 有关。
  \hfill\xstarfull{3}\yiju{E6 教材 5.4；E2 slide5}\nandu{中}

  \item 流体是一种在任一剪切力的作用下不能保持静止的物质。\hfill\xstarfull{2}\yiju{概念题素材；E6 教材 1.1}\nandu{易}

  \item 均质不可压缩流体是指密度为常数的流体。\hfill\xstarfull{3}\yiju{E4 2015 填空 1（由 $\rho=851$ kg/m$^3$ 求重度与动力黏度）}\nandu{易}

  \item 气体的粘度随着温度的升高而减小。\hfill\xstarfull{5}\yiju{E4 2026 计算 1（牛顿流体斜板静压强）；E5 题库选择 2（汽油＝牛顿流体）}\nandu{中}

  \item 在圆管流中，层流的断面流速分布符合抛物线分布规律。\hfill\xstarfull{5}\yiju{E4 2022 单选 10（管径减半压强损失 16 倍）；E4 2015 填空（层流平均速度＝0.5 倍最大速度）}\nandu{易}

  \item 当水平放置的管道中水流为恒定均匀流时，其动水压强沿程减小。\hfill\xstarfull{5}\yiju{E1 slide16"计算题：伯努利方程"；E1 slide33 与 E2 slide4"第四章＝考试计算题重点"}\nandu{中}
\end{ti}

\begin{tishi}
\textbf{第 4 题：}原卷此行文字 OCR 为"恒定流动中，流线和迹线总是重合的"；同类题库中常见
"恒定流动中流线与迹线不重合"之说，二者结论相反。此处\textbf{按原卷 OCR 文字转写}，
请对照原卷影印核实。
\end{tishi}

\section{三、填空题（共 30 分，每空 3 分）}

\begin{ti}
  \item 研究流体运动的方法有拉格朗日法和 \underline{\hspace{3cm}}。\hfill\xstarfull{5}\yiju{E4 各年选择/填空（质点导数与加速度）；E1 slide33"公式较多，重点记忆"}\nandu{易}

  \item 流体微团的运动形式有平移、\underline{\hspace{2.5cm}} 和 \underline{\hspace{2.5cm}}。
  \hfill\xstarfull{4}\yiju{E5 题库选择 13（涡量为零）；E6 教材 3.4}\nandu{中}

  \item 任意空间点上的运动参数都不随时间变化的流动称为 \underline{\hspace{3cm}}。
  \hfill\xstarfull{3}\yiju{E1 slide16 简答（均匀流的特点）}\nandu{易}

  \item 雷诺实验揭示了液体存在 \underline{\hspace{2.5cm}} 和 \underline{\hspace{2.5cm}} 两种流态。
  \hfill\xstarfull{5}\yiju{E4 2022 单选 9（两圆管流态判别）；E4 2026 简答 3（层流与雷诺数关系）}\nandu{易}

  \item 层流的沿程水头损失与管截面上的平均流速的 \underline{\hspace{2cm}} 次方成正比。
  \hfill\xstarfull{5}\yiju{E4 2022 单选 3；E5 题库选择 10（圆管与方管沿程损失比 0.886）}\nandu{中}

  \item 不可压缩流体的空间三维的连续性微分方程是 \underline{\hspace{4cm}}。
  \hfill\xstarfull{5}\yiju{E4 2022 单选 4；E5 题库}\nandu{中}

  \item 如图 2 所示，圆管突然扩大的水头损失 $h_j$ 可表示为 \underline{\hspace{4cm}}。
  \hfill\xstarfull{4}\yiju{E4 2015 计算 3（突然缩小水头损失）；E1 slide16}\nandu{中}
  \tuyi{原卷图 2：一段圆管由粗管突然扩大为细管（沿流向管径由 $d_1$ 突变为 $d_2$），
  在突变断面处产生局部水头损失 $h_j$；原卷图中标注有流速 $v$，未给出具体数值。
  本题只要求写出 $h_j$ 的表达式（常用 Borda--Carnot 公式 $h_j=\zeta\dfrac{v^{2}}{2g}$，
  $\zeta$ 按所取断面确定）。（原卷影印见本册附录）}

  \item 如图 3 所示，一高为 $H$、直径为 $D$ 的旋转圆柱容器，容器中盛水，当水正好与容器中心
  触底、顶部与容器同高时，容器的旋转角速度等于 \underline{\hspace{4cm}}。
  \hfill\xstarfull{3}\yiju{E6 教材 2.4；E5 题库}\nandu{难}
  \tuyi{原卷图 3：一圆柱形容器（高 $H$、直径 $D$）绕自身铅垂中心轴以角速度 $\omega$ 旋转，
  水面形成旋转抛物面；题设水面最低点（抛物面顶点）恰在容器底面中心，最高点与容器顶缘齐平。
  原卷图中另有半径标记"$R$"，但题干未给出 $R$ 的数值，故答案应以已知量 $H$、$D$ 表示
  （$R=D/2$）。（原卷影印见本册附录）}
\end{ti}

\section{四、计算题（共 85 分）}

\noindent （1、2、3、4、5、6、7 题必做，8、9、10、11 题任选做两道。）

\begin{ti}
  \item 如图 4 所示，一底面积为 $40\times45\,\mathrm{cm}$、高为 $1\,\mathrm{cm}$ 的木板，
  质量为 $5\,\mathrm{kg}$，沿着涂有润滑油的斜面等速向下运动。已知 $v=1\,\mathrm{m/s}$，
  $\delta=1\,\mathrm{mm}$，求润滑油的动力粘滞系数。（9 分）\hfill\xstarfull{5}\yiju{E4 2026 计算 1（牛顿流体斜板静压强）；E5 题库选择 2（汽油＝牛顿流体）}\nandu{中}
  \tuyi{原卷图 4：一木板置于倾角为 $\theta$ 的斜面上，板面与斜面之间充满厚 $\delta$ 的润滑油层，
  木板沿斜面向下以 $v=1\,\mathrm{m/s}$ 等速运动；板沿斜面方向长为 $l$，宽度方向待考证
  （需与底面积 $40\times45\,\mathrm{cm}$ 及斜面倾角配合才能算出粘滞系数）。
  （原卷影印见本册附录）}

  \item 如图 5 所示，有两个水池，其底部以一水管连通。在恒定的水面差 $H$ 的作用下，
  水从左水池流入右水池。水管直径 $d=500\,\mathrm{mm}$，管总长 $100\,\mathrm{m}$，直角进口，
  弯管呈 $90^\circ$，管中流量为 $0.5\,\mathrm{m^2/s}$。已知弯管 $\zeta_{\text{弯}}=0.1455$，
  $\zeta_{\text{进}}=0.81$，$\lambda=0.021$，求两水池水面的高差 $H$。（10 分）
  \hfill\xstarfull{5}\yiju{E5 题库计算 2（水箱垂直管道出流 $Q$ 与管长 $l$）；E1 slide33}\nandu{中}
  \tuyi{原卷图 5：左、右两个敞口水池，底部用一根水管连通，管路上有一个 $90^\circ$ 弯管，
  水从左侧水池经该管流入右侧水池；两池水面高差为 $H$。（原卷影印见本册附录）}

  \item 如图 6 所示，两圆筒用管子连接，内充水银。第一个圆筒直径 $d_1=45\,\mathrm{cm}$，
  活塞上受力 $F_1=3197\,\mathrm{N}$，密封气体的计示压强 $p_e=9810\,\mathrm{Pa}$；
  第二个圆筒直径 $d_2=30\,\mathrm{cm}$，活塞上受力 $F_2=4945.5\,\mathrm{N}$，开口通大气。
  若不计活塞质量，求平衡状态时两活塞的高度差 $h$。
  （已知水银密度 $\rho=13600\,\mathrm{kg/m^3}$）（8 分）\hfill\xstarfull{4}\yiju{E4 2018 选择 3（U 形水银测压计）；E4 2015 填空（U 形压差计）}\nandu{中}
  \tuyi{原卷图 6：两个竖直圆筒下端用一根管子连通，管内及两筒下部充水银；
  左筒（直径 $d_1$）活塞上作用向下的力 $F_1$，其上部气体计示压强为 $p_e$；
  右筒（直径 $d_2$）活塞上作用向下的力 $F_2$，筒口通大气；
  两活塞（水银面）的高差为 $h$。（原卷影印见本册附录）}

  \item 如图 7 所示，水由喷嘴射出，已知流量 $Q=0.4\,\mathrm{m^3/s}$，主管直径 $D=0.4\,\mathrm{m}$，
  喷口直径 $d=0.1\,\mathrm{m}$，水头损失不计，求水平方向上水流作用在喷嘴上的力。（10 分）
  \hfill\xstarfull{5}\yiju{E4 2015 计算（水平弯管求 $F_x,F_y$）；E4 2026 计算 3（教材习题 4.12）；E1 slide16}\nandu{难}
  \tuyi{原卷图 7：水平放置的渐缩管（喷嘴），进口主管直径 $D$、出口喷口直径 $d$，
  水沿水平方向由喷口射入大气；取喷嘴内壁为控制体表面，用连续性方程与动量方程求轴向力。
  （原卷影印见本册附录）}

  \item 如图 8 所示，水在变直径竖管中流动，已知粗管直径 $d_1=300\,\mathrm{mm}$，
  流速 $v_1=6\,\mathrm{m/s}$。为使两断面的压力表读值相同，试求细管直径
  （水头损失不计）。（10 分）\hfill\xstarfull{5}\yiju{E1 slide16"计算题：伯努利方程"；E1 slide33 与 E2 slide4"第四章＝考试计算题重点"}\nandu{难}
  \tuyi{原卷图 8：竖直放置的变直径管段，下部为粗管断面 1--1（直径 $d_1=300\,\mathrm{mm}$，
  流速 $v_1=6\,\mathrm{m/s}$），上部为细管断面 2--2，两断面各接压力表；
  两断面之间的高差记为 $h$，题目要求使两压力表读值相同（$p_1=p_2$）时的细管直径 $d_2$。
  （原卷影印见本册附录）}

  \item 如图 9，已知 $h_1=600\,\mathrm{mm}$，$h_2=250\,\mathrm{mm}$，$h_3=200\,\mathrm{mm}$，
  $h_4=300\,\mathrm{mm}$，$h_5=500\,\mathrm{mm}$，$\rho_1=1000\,\mathrm{kg/m^3}$，
  $\rho_2=800\,\mathrm{kg/m^3}$，$\rho_3=13598\,\mathrm{kg/m^3}$，求 $A$、$B$ 两点的压强差。（8 分）
  \hfill\xstarfull{4}\yiju{E4 2018 选择 3（U 形水银测压计）；E4 2015 填空（U 形压差计）}\nandu{难}
  \tuyi{原卷图 9：由多根竖管与 U 管串联而成的复式压差计，$A$、$B$ 为被测两点；
  管中交替充以密度为 $\rho_1$、$\rho_2$、$\rho_3$ 的三种流体，
  各段液柱高度依次为 $h_1$、$h_2$、$h_3$、$h_4$、$h_5$；
  各段上下连接顺序与左右走向请对照原卷影印。（原卷影印见本册附录）}

  \item 平面不可压缩流体的速度势函数 $\varphi=x^{2}-y^{2}-x$，
  求流场上 $A(-1,\,-1)$ 及 $B(2,\,2)$ 点处的速度值及流函数值。（10 分）
  \hfill\xstarfull{5}\yiju{E4 2015 计算（流函数求速度与流线）；E4 2026 计算 4（证明势函数存在性）}\nandu{中}

  \item 如图 10 所示，平面上有一对等强度为 $\Gamma$（$\Gamma>0$）的点涡，其方向相反，
  分别位于 $(0,\,h)$、$(0,\,-h)$ 两固定点处，同时平面上有一无穷远平行于 $x$ 轴的来流 $v_\infty$，
  试求合成速度在原点的值。（10 分）\hfill\xstarfull{2}\yiju{概念题；E6 教材 7.8}\nandu{难}
  \tuyi{原卷图 10：平面直角坐标系 $Oxy$ 中，$y$ 轴上关于原点对称的两点 $(0,h)$、$(0,-h)$
  各置一个强度均为 $\Gamma$、旋转方向相反的点涡；另有沿 $x$ 轴正方向的无穷远均匀来流
  $v_\infty$。（原卷影印见本册附录）}

  \item 平底船的底面可视为宽 $b=10\,\mathrm{m}$、长 $L=50\,\mathrm{m}$ 的平板，船速
  $v_0=4\,\mathrm{m/s}$，水的运动粘度 $\nu=10^{-6}\,\mathrm{m^2/s}$，
  如果平板边界层转换临界雷诺数 $\Rey_{xcr}=5\times10^{5}$，试求克服边界层阻力所需的功率。（10 分）
  \hfill\xstarfull{2}\yiju{E4 2015 计算 9（平板摩擦阻力比 $F_1/F_2$）}\nandu{难}

  \item 过热蒸汽（$k=1.33$，$R=462\,\mathrm{J/(kg\cdot K)}$）的温度为 $430\,^\circ\mathrm{C}$，
  压强为 $5\times10^{5}\,\mathrm{Pa}$，速度为 $525\,\mathrm{m/s}$，
  求过热蒸汽的滞止参数 $T_0$、$p_0$ 和 $\rho_0$。（10 分）\hfill\xstarmark{X2}\nandu{难}

  \item 试确定某梯形断面均匀流渠道的水力最佳断面尺寸。已知边坡系数 $m=1.5$，
  糙率 $n=0.025$，底坡 $i=0.002$，流量 $Q=3.0\,\mathrm{m^3/s}$。（10 分）
  \hfill\xstarmark{X1}\nandu{难}
\end{ti}

\begin{tishi}
\textbf{第 1 题（计算）：}原卷数值 OCR 不清——本题给出的"底面积 $40\times45\,\mathrm{cm}$、
高为 $1\,\mathrm{cm}$"与"$\delta=1\,\mathrm{mm}$"互相矛盾（油膜厚度应为 $\delta$，
而不是木板"高 $1\,\mathrm{cm}$"），且题干未抄出斜面倾角与沿流向的板长。
以上数值\textbf{一律照 OCR 保留}，未作修改，请对照原卷影印核实后再作答。

\textbf{第 2 题（计算）：}原卷 OCR 将流量写作"$0.5\,\mathrm{m^2/s}$"，流量单位应为
$\mathrm{m^3/s}$，此处已按单位改正；\textbf{数值 $0.5$ 未改}。

\textbf{第 3 题（计算）：}原卷 OCR 将水银密度写作"$13600\,\mathrm{kg/m^2}$"，密度单位应为
$\mathrm{kg/m^3}$，此处已按单位改正；\textbf{数值 $13600$ 未改}。

\textbf{第 4 题（计算）：}原卷 OCR 将流量写作"$0.4\,\mathrm{m^2/s}$"，应为 $\mathrm{m^3/s}$，
此处已按单位改正；\textbf{数值 $0.4$ 未改}。

\textbf{第 5 题（计算）：}原卷数值 OCR 不清——题干与图中应有的\textbf{两断面高差 $h$}
（或细管断面位置）未抄录出来，而它是解本题的关键；照现有条件只能给出
$d_2=d_1\sqrt{v_1/(v_1-2gh)}$。上式照原卷保留、未增补任何数值，请对照原卷影印核实。

\textbf{第 6 题（计算）：}原卷数值 OCR 不清——图中各段竖管与 $h_1,\dots,h_5$ 的对应关系、
各段液体的连接顺序 OCR 未给出，此处按题干列出的数值照录，请对照原卷影印核实。

\textbf{第 8 题（计算）：}原卷数值 OCR 不清——来流速度 $v_\infty$ 只以符号出现在题干中，
原卷未抄出其数值，故答案只能写成含 $v_\infty$、$\Gamma$、$h$ 的表达式。

\textbf{第 9 题（计算）：}原卷数值 OCR 不清——临界雷诺数 OCR 作"$5\times10^{2}$"，
按平板边界层转捩的通用取值应为 $\Rey_{xcr}=5\times10^{5}$；此处已按通用取值改正。
另，阻力计算所取的\textbf{湿面积}（单面还是双面）原卷未明确，请对照原卷影印核实。

\textbf{第 10 题（计算）：}原卷数值 OCR 不清——压强一项 OCR 作"$5\times10^{5}\,\mathrm{Pa}$"
（原文指数位字形模糊），此处照该读法转写、\textbf{未改动数值}，请对照原卷影印核实。
本题属\textbf{气体动力学（一元可压缩流动）}内容，为旧大纲/复试范围，现行 818 大纲不要求。

\textbf{第 11 题（计算）：}原卷数值 OCR 不清——流量 OCR 作"$Q=3.0\,\mathrm{m^3/s}$"。
按 $Q=3.0\,\mathrm{m^3/s}$、$m=1.5$、$n=0.025$、$i=0.002$ 用水力最佳断面条件试算，
所得断面水深与底宽偏小（明渠中该组合的量级偏小），故流量的数量级可能 OCR 有误；
此处\textbf{照原卷数值保留、未改}，请对照原卷影印核实。本题为\textbf{旧大纲（2018 年及以前）}
内容，现行 818 大纲不要求明渠流与堰流。
\end{tishi}
